\chapter{Il Web e il protocollo HTTP}

% ══════════════════════════════════════════════════════════════
\section{Dal Web a HTTP}
% ══════════════════════════════════════════════════════════════

Fino ai primi anni '90 Internet era usata soprattutto da ricercatori, docenti e studenti
universitari: per collegarsi a host remoti (Telnet), trasferire file (FTP), leggere e inviare
news e posta elettronica. Applicazioni utilissime, ma che lasciavano Internet sconosciuta fuori
dal mondo accademico. Nei primi anni '90 arrivò il \textbf{World Wide Web}, la prima applicazione
Internet a conquistare il grande pubblico.

\dfn{World Wide Web}{
    Il \textbf{World Wide Web} (WWW, o semplicemente Web) è una collezione di informazioni
    (documenti, immagini, video, audio, \dots) accessibili attraverso Internet secondo uno specifico
    protocollo, l'\textbf{HTTP} (\emph{HyperText Transfer Protocol}).
}

La comunicazione HTTP è tipicamente \textbf{client-server}:

\begin{itemize}
    \item un \textbf{programma client} traduce le richieste dell'utente in messaggi HTTP: è implementato
          nei \textbf{browser} (Chrome, Firefox, Edge, Safari, \dots);
    \item un \textbf{programma server} esegue la richiesta HTTP e restituisce una risposta HTTP.
\end{itemize}

\begin{figure}[H]
\centering
\begin{tikzpicture}[every node/.style={font=\footnotesize\sffamily}]
  \node (srv) at (0,0) {\icon[0.9cm]{server}};
  \node[below=-1pt of srv, align=center] {server Web\\(per esempio Apache)};
  \node (c1) at (-5,1.3) {\icon[1.1cm]{pc}}; \node[below=-2pt of c1] {PC con Firefox};
  \node (c2) at (5,1.3) {\icon[1.2cm]{laptop}}; \node[below=-2pt of c2] {laptop con Safari};
  \node (c3) at (4.2,-1.6) {\icon[0.8cm]{smartphone}}; \node[below=-2pt of c3] {smartphone con Chrome};
  \foreach \c/\sx in {c1/-0.5, c2/0.5, c3/0.4} {
    \draw[msg, netblue, shorten <=14pt, shorten >=14pt] ([yshift=4pt]\c.center) -- ([yshift=4pt, xshift=\sx cm]srv.center);
    \draw[msg, netgreen, shorten <=14pt, shorten >=14pt] ([yshift=-4pt, xshift=\sx cm]srv.center) -- ([yshift=-4pt]\c.center);
  }
  \draw[msg, netblue] (-6.2,-2.2) -- ++(0.8,0) node[right, text=netblue] {richiesta HTTP};
  \draw[msg, netgreen] (-6.2,-2.7) -- ++(0.8,0) node[right, text=netgreen] {risposta HTTP};
\end{tikzpicture}
\caption{Client diversi, con browser diversi, parlano con lo stesso server perché condividono il protocollo HTTP.}
\end{figure}

HTTP definisce principalmente due cose: la \textbf{struttura dei messaggi} e \textbf{come client e
server se li scambiano}.

\warning{Cosa HTTP non definisce}{
    HTTP non stabilisce come vadano implementati client e server, né come i contenuti vadano
    \textbf{visualizzati}. Il modo in cui una pagina appare dipende dal browser: due browser possono
    mostrare la stessa pagina in modi un po' diversi (uno aspetta che sia tutta caricata, l'altro la
    mostra pezzo per pezzo).
}

Il Web e i suoi protocolli sono la piattaforma di applicazioni come YouTube, la posta via Web
(Gmail) e gran parte delle app per dispositivi mobili, social network compresi.

\begin{figure}[H]
\centering
\begin{tikzpicture}[every node/.style={font=\footnotesize\sffamily}]
  \node (c) at (0,0) {\icon[1cm]{laptop}};
  \node[lay app, minimum width=2.4cm] (ca) at (0,-1.1) {browser};
  \node[lay tra, minimum width=2.4cm] (ct) at (0,-1.72) {};
  \node[lbl] at (0,-2.3) {[\dots] resto della pila};
  \node (s) at (7,0) {\icon[0.7cm]{server}};
  \node[lay app, minimum width=2.4cm] (sa) at (7,-1.1) {server Web};
  \node[lay tra, minimum width=2.4cm] (st) at (7,-1.72) {};
  \node[lbl] at (7,-2.3) {[\dots] resto della pila};
  \draw[{Stealth}-{Stealth}, thick, netred] (ca) -- node[above, text=netred] {HTTP} (sa);
  \draw[{Stealth}-{Stealth}, thick, netred] (ct) -- node[below, text=netred] {TCP (o UDP)} (st);
  \node[anchor=west] at (9,-1.1) {livello applicativo};
  \node[anchor=west] at (9,-1.72) {livello di trasporto};
\end{tikzpicture}
\caption{HTTP è il protocollo tra browser e server Web; sotto, il trasporto è affidato a TCP (o, con HTTP/3, a UDP).}
\end{figure}

% ══════════════════════════════════════════════════════════════
\section{Le URL}
% ══════════════════════════════════════════════════════════════

\dfn{Risorsa e URL}{
    Le informazioni sul Web si chiamano \textbf{risorse} (o \textbf{oggetti}) e sono identificate da un
    \textbf{URL} (\emph{Uniform Resource Locator}), una stringa con la forma generale:
    \begin{center}
    \texttt{[protocollo]://[userinfo@][host][:porta][/percorso][?query][\#frammento]}
    \end{center}
}

\begin{figure}[H]
\centering
\begin{tikzpicture}[every node/.style={font=\ttfamily\large, inner sep=2pt}]
  \node[fill=lapp!60] (p) at (0,0) {https};
  \node[right=0pt of p] (s1) {://};
  \node[fill=lnet!60, right=0pt of s1] (h) {en.wikipedia.org};
  \node[fill=ltra!60, right=0pt of h] (pt) {:443};
  \node[fill=llink!80, right=0pt of pt] (pa) {/w/index.php};
  \node[fill=netred!20, right=0pt of pa] (q) {?title=SSC\_Napoli};
  \node[fill=lphy, right=0pt of q] (f) {\#Storia};
  \foreach \n/\t/\d in {p/protocollo/0.5, h/{host (nome o indirizzo IP)}/0.5, pt/porta/1.2, pa/percorso/0.5, q/query/1.2, f/frammento/0.5}
     \draw[{Stealth}-, thin] (\n.south) -- ++(0,-\d) node[below, font=\small\sffamily] {\t};
\end{tikzpicture}
\caption{Anatomia di un URL.}
\end{figure}

\begin{itemize}
    \item \texttt{[protocollo]}: il protocollo con cui accedere alla risorsa (HTTP, HTTPS, FTP, \dots).
    \item \texttt{[userinfo]} (opzionale): informazioni sull'utente, come \texttt{nome:password@}. È
          \textbf{deprecato} per ovvi motivi di sicurezza.
    \item \texttt{[host]}: il nome o l'indirizzo IP del server.
    \item \texttt{[porta]} (opzionale): la porta da usare; di solito si ricava dal protocollo (80 per HTTP,
          443 per HTTPS).
    \item \texttt{[percorso]}: dove si trova la risorsa dentro il server (\texttt{/percorso/della/risorsa}).
    \item \texttt{[query]}: preceduta da \texttt{?}, specifica eventuali richieste (parametri).
    \item \texttt{[frammento]}: preceduto da \texttt{\#}, identifica un elemento dentro la risorsa (una
          sezione, un modulo).
\end{itemize}

\begin{esempio}{Due URL scomposti}
\texttt{http://www.someSchool.edu/someDepartment/picture1.gif}
\begin{itemize}
    \item protocollo: \texttt{http}; host: \texttt{www.someSchool.edu}; percorso: \texttt{/someDepartment/picture1.gif}.
\end{itemize}
\texttt{https://en.wikipedia.org/w/index.php?title=SSC\_Napoli}
\begin{itemize}
    \item protocollo: \texttt{https}; host: \texttt{en.wikipedia.org}; percorso: \texttt{/w/index.php};
          query: \texttt{?title=SSC\_Napoli}.
\end{itemize}
In questo caso la query produce la stessa pagina di \texttt{https://en.wikipedia.org/wiki/SSC\_Napoli}.
\end{esempio}

\subsection{Pagine e oggetti}

La maggior parte delle pagine Web è formata da un \textbf{file HTML di base} (\emph{HyperText Markup
Language}) e da vari \textbf{riferimenti} ad altri oggetti (immagini, video, script, fogli di stile),
ciascuno con il proprio URL.

\begin{esempio}{Quanti oggetti ha una pagina}
Una pagina con un testo HTML e cinque immagini JPEG ha \textbf{sei oggetti}: il file HTML di base più le
cinque immagini. Il browser riceve prima l'HTML, vi trova i riferimenti alle immagini e le chiede una a una.
\end{esempio}

Una pagina, quindi, non è un trasferimento ma \textbf{diversi} trasferimenti: ed è proprio il modo di
organizzarli il tema del resto del capitolo.

\subsection{Prima del primo messaggio HTTP}

Tra il momento in cui si preme Invio e quello in cui parte la prima richiesta succedono già parecchie cose:

\begin{enumerate}
    \item il browser \textbf{scompone l'URL}: protocollo, host, porta, percorso;
    \item l'host è un \textbf{nome}, e verso un nome non si può aprire nessuna connessione: il \textbf{DNS}
          fornisce l'indirizzo IP corrispondente (lo vedremo nel capitolo dedicato);
    \item solo a questo punto si apre una \textbf{connessione TCP} verso quell'indirizzo, sulla porta 80 o 443;
    \item su quella connessione viaggia la prima \textbf{richiesta HTTP}, con il percorso;
    \item il browser legge l'HTML, trova gli oggetti a cui rimanda e ricomincia.
\end{enumerate}

\info{Il tempo che non contiamo}{
    Anche la risoluzione del nome costa tempo, ma nei conti in RTT che seguono non viene mai inclusa:
    si parte sempre da quando l'indirizzo del server è già noto.
}

% ══════════════════════════════════════════════════════════════
\section{Come funziona HTTP}
% ══════════════════════════════════════════════════════════════

HTTP definisce come i client Web chiedono oggetti ai server Web e come i server li trasferiscono.
Quando l'utente chiede un oggetto (per esempio cliccando un link), il browser invia messaggi di
\textbf{richiesta HTTP} per gli oggetti della pagina; il server risponde con messaggi di
\textbf{risposta HTTP} che contengono gli oggetti.

HTTP evolve continuamente per bilanciare affidabilità e prestazioni:

\begin{figure}[H]
\centering
\begin{tikzpicture}[every node/.style={font=\footnotesize\sffamily}]
  \foreach \x/\v/\y/\c in {0/{HTTP/1.0}/1996/lnet, 3.3/{HTTP/1.1}/1997/lnet, 6.6/{HTTP/2}/2015/ltra, 9.9/{HTTP/3}/2022/lapp} {
    \node[signal, signal to=east, fill=\c, draw=black!50, minimum width=3.1cm, minimum height=0.9cm, align=center] at (\x,0) {\bfseries \v\\\y};
  }
  \node[lbl, text width=3cm] at (0,-0.9) {connessioni non persistenti};
  \node[lbl, text width=3cm] at (3.3,-0.9) {connessioni persistenti};
  \node[lbl, text width=3cm] at (6.6,-0.9) {richieste intercalate e priorità};
  \node[lbl, text width=3cm] at (9.9,-0.9) {su UDP con QUIC};
\end{tikzpicture}
\caption{Le versioni di HTTP.}
\end{figure}

\subsection{Il protocollo di trasporto sottostante}

HTTP usa principalmente \textbf{TCP} come protocollo di trasporto; dal 2022, con \textbf{HTTP/3}, può
usare anche \textbf{UDP} (circa il 30\% del traffico HTTP viaggia oggi su UDP).

\begin{enumerate}
    \item Il client HTTP \textbf{apre una connessione} TCP con il server.
    \item Stabilita la connessione, browser e server comunicano attraverso i loro \textbf{socket}: il client
          scrive richieste nel proprio socket e legge risposte dal proprio socket; il server fa il contrario.
\end{enumerate}

Perché storicamente TCP? Perché offre servizi utilissimi, soprattutto il \textbf{trasferimento
affidabile}: garantisce che richieste e risposte arrivino integre a destinazione (se possibile). Ma i
servizi di TCP rallentano la comunicazione: per questo HTTP/3 usa UDP con sopra il protocollo
\textbf{QUIC} (\emph{Quick UDP Internet Connections}), che reimplementa l'affidabilità sopra UDP ed è
stimato essere circa il 30\% più veloce.

\info{Il vantaggio della stratificazione, visto da HTTP}{
    Qui si vede uno dei grandi vantaggi dell'architettura a strati: le applicazioni HTTP non devono
    preoccuparsi di raggiungibilità, perdite di dati, ritrasmissioni. Possono essere più semplici, sia
    logicamente sia computazionalmente, delegando gran parte del lavoro agli strati sottostanti.
}

% ══════════════════════════════════════════════════════════════
\section{HTTP è stateless}
% ══════════════════════════════════════════════════════════════

La semplicità è fondamentale per un server che riceve moltissime richieste al secondo (in media
un server gestisce circa 1000 richieste HTTP al secondo, e l'infrastruttura di un grande motore di
ricerca centinaia di migliaia di interrogazioni al secondo).

\dfn{Protocollo stateless}{
    HTTP è \textbf{senza stato} (\emph{stateless}): il server \textbf{non conserva alcuna informazione}
    sulle interazioni passate con un certo client.
}

\begin{esempio}{La stessa richiesta due volte}
Se un client chiede lo stesso oggetto due volte di seguito, il server non considera la seconda
richiesta ridondante: rimanda semplicemente l'oggetto, perché ha completamente dimenticato la prima.
\end{esempio}

Non tutte le applicazioni sono stateless: molte sono \textbf{stateful} (con stato). Vedremo come HTTP
riesca comunque a ``ricordarsi'' degli utenti usando i \textbf{cookie}.

% ══════════════════════════════════════════════════════════════
\section{Connessioni persistenti e non persistenti}
% ══════════════════════════════════════════════════════════════

Client e server possono comunicare a lungo, scambiandosi molte coppie richiesta--risposta (una di
seguito all'altra, a intervalli regolari o in modo intermittente). Con un trasporto orientato alla
connessione (TCP o QUIC) ci sono due possibilità.

\dfn{Connessioni persistenti e non persistenti}{
    \begin{itemize}
      \item \textbf{Connessione non persistente}: per \textbf{ogni} coppia richiesta--risposta si apre una
            \textbf{nuova} connessione, che viene chiusa subito dopo.
      \item \textbf{Connessione persistente}: tutte le richieste e le risposte viaggiano sulla
            \textbf{stessa} connessione.
    \end{itemize}
}

Le prime versioni di HTTP (HTTP/1.0) usavano per default connessioni non persistenti; oggi il default
sono le connessioni persistenti, ma client e server possono essere configurati per usare quelle non
persistenti.

\subsection{Esempio di connessione non persistente}

Chiediamo una pagina fatta di un file HTML di base e \textbf{10 immagini JPEG}, tutte sullo stesso
server, all'URL \url{http://www.someschool.edu/someDepartment/home.index}.

\begin{enumerate}
    \item Il client HTTP \textbf{apre una connessione TCP} verso \texttt{www.someschool.edu} sulla porta
          \textbf{80}: ci sono ora due socket, uno sul client e uno sul server.
    \item Il client invia nel proprio socket una \textbf{richiesta HTTP} che contiene il percorso
          \texttt{/someDepartment/home.index}.
    \item Il server riceve la richiesta, recupera l'oggetto (da RAM o disco), lo incapsula in una
          \textbf{risposta HTTP} e la invia nel proprio socket.
    \item Il server dice a TCP di \textbf{chiudere} la connessione; TCP la chiuderà davvero quando sarà
          sicuro che il client ha ricevuto la risposta integra.
    \item Il client riceve la risposta e la connessione termina. Il messaggio dice che l'oggetto è un file
          HTML: il client lo estrae, lo esamina e trova i riferimenti ai 10 JPEG.
    \item I passi 1--4 vengono ripetuti \textbf{per ciascuna delle 10 immagini}.
\end{enumerate}

\begin{figure}[H]
\centering
\begin{tikzpicture}[every node/.style={font=\footnotesize\sffamily}]
  \node (c) at (0,0) {\icon[1.1cm]{pc}}; \node[below=-2pt of c] {client};
  \node (s) at (8,0) {\icon[0.8cm]{server}}; \node[below=-2pt of s] {server};
  \foreach \i/\y/\t in {1/1.5/HTML, 2/0.9/JPEG 1, 3/0.3/JPEG 2, 11/-1.2/JPEG 10} {
    \draw[{Stealth}-{Stealth}, thick, netblue] (1,\y) -- node[above=-1pt, fill=white, inner sep=1pt] {connessione TCP n.~\i: \t} (7,\y);
  }
  \node at (4,-0.45) {$\vdots$};
\end{tikzpicture}
\caption{Con connessioni non persistenti servono 11 connessioni TCP (e 11 coppie di socket) per una pagina di 11 oggetti.}
\end{figure}

Poiché la connessione non sopravvive tra un oggetto e l'altro, servono \textbf{11 connessioni TCP} per
trasferire la pagina, e gli elementi possono arrivare in momenti diversi.

\subsection{Connessioni in sequenza e in parallelo}

Le connessioni possono essere aperte \textbf{in sequenza} o \textbf{in parallelo}. Nell'esempio, una volta
ricevuto l'HTML e scoperto che servono 10 immagini, il browser può scaricarle tutte
contemporaneamente.

\begin{itemize}
    \item I browser moderni permettono di configurare il grado di parallelismo: di solito aprono \textbf{da 5
          a 10 connessioni TCP parallele}, ciascuna con una sola transazione richiesta--risposta.
    \item Le connessioni parallele \textbf{accorciano il tempo di risposta} ma \textbf{aumentano il costo
          computazionale}.
    \item Anche con il parallelismo, aprire tante connessioni TCP ha un costo in tempo non trascurabile.
\end{itemize}

% ══════════════════════════════════════════════════════════════
\section{Il costo in RTT}
% ══════════════════════════════════════════════════════════════

È difficile stimare con precisione i tempi su Internet. Si può però misurare il costo di una richiesta
in \textbf{RTT}.

\dfn{Round-Trip Time (RTT)}{
    L'\textbf{RTT} è il tempo impiegato da un \textbf{piccolo pacchetto} per andare dal client al server e
    \textbf{tornare} al client.
}

\subsection{Il three-way handshake}

Con un trasporto orientato alla connessione come TCP, prima di scambiare dati si esegue il
\textbf{three-way handshake} (stretta di mano a tre vie):

\begin{figure}[H]
\centering
\begin{tikzpicture}[yscale=0.85, every node/.style={font=\footnotesize\sffamily}]
  \timeline{c}{0}{Client}{5.2}
  \timeline{s}{6}{Server}{5.2}
  \draw[msg, netblue] (0,-0.4) -- node[above, sloped] {1. ``vorrei aprire una connessione'' (SYN)} (6,-1.4);
  \draw[msg, netblue] (6,-1.6) -- node[above, sloped] {2. ``ricevuto, sono pronto'' (SYN-ACK)} (0,-2.6);
  \draw[msg, netgreen] (0,-2.8) -- node[above, sloped] {3. ``ricevuto'' + richiesta HTTP (ACK)} (6,-3.8);
  \draw[msg, netgreen, very thick] (6,-4.0) -- node[above, sloped] {risposta: il file HTML} (0,-5.0);
  \draw[decorate, decoration={brace, amplitude=5pt, mirror}] (-0.2,-0.4) -- (-0.2,-2.6) node[midway, left=6pt] {1 RTT};
  \draw[decorate, decoration={brace, amplitude=5pt, mirror}] (-0.2,-2.8) -- (-0.2,-5.0) node[midway, left=6pt] {1 RTT};
  \draw[decorate, decoration={brace, amplitude=5pt}] (6.2,-4.0) -- (6.2,-4.4) node[midway, right=6pt, text width=2.6cm, align=left] {tempo di trasmissione del file};
  \node[right, text=netgreen!60!black] at (6.3,-3.1) {connessione stabilita};
\end{tikzpicture}
\caption{Richiesta di un file HTML con connessione non persistente: circa 2 RTT più il tempo di trasmissione del file.}
\end{figure}

\begin{enumerate}
    \item Il client invia un piccolo segmento TCP per segnalare che vuole aprire una connessione.
    \item Il server conferma e risponde con un piccolo segmento TCP.
    \item Il client conferma a sua volta.
\end{enumerate}

Dal punto di vista degli RTT, i primi due passi dell'handshake richiedono \textbf{1 RTT}. Un
handshake analogo avviene anche quando la connessione viene chiusa.

\subsection{Il conto totale}

\begin{itemize}
    \item In HTTP il client invia la \textbf{richiesta} insieme al terzo passo dell'handshake (la conferma).
    \item Quando la richiesta arriva, il server invia il file HTML: la coppia richiesta--risposta costa
          \textbf{un altro RTT}.
\end{itemize}

\dfn{Tempo di risposta con connessione non persistente}{
    $$T_{\text{risposta}} \approx \underbrace{1\ \text{RTT}}_{\text{apertura TCP}} + \underbrace{1\ \text{RTT}}_{\text{richiesta e risposta}} + t_{\text{trasmissione del file}} = 2\ \text{RTT} + t_{\text{trasmissione}}$$
}

\begin{esempio}{La stessa pagina, contata in quattro modi}
La pagina con l'HTML e 10 immagini JPEG (11 oggetti), senza perdite e con tempi di trasmissione
trascurabili, come nelle slide:

\begin{center}
\renewcommand{\arraystretch}{1.3}
\begin{tabular}{p{6.2cm}p{4.6cm}r}
\toprule
\textbf{Modalità} & \textbf{Conto} & \textbf{RTT} \\
\midrule
Non persistente, un oggetto alla volta & $2 \times 11$ & \textbf{22} \\
Non persistente, 5 connessioni in parallelo & 2 per l'HTML, poi due turni da 2 per le 10 immagini & \textbf{6} \\
Persistente, senza pipelining & 1 per aprire, poi 1 per ogni oggetto & \textbf{12} \\
Persistente con pipelining & 1 per aprire, 1 per l'HTML, 1 per le 10 immagini insieme & \textbf{3} \\
\bottomrule
\end{tabular}
\end{center}

Con RTT $= 50$ ms si va da $22 \times 50 = 1{,}1$ s a $3 \times 50 = 0{,}15$ s. Stessa pagina, stesso server,
stessa distanza: tutto ciò che sta fra 22 e 3 dipende solo da \textbf{come vengono gestite le connessioni}.
\end{esempio}

% ══════════════════════════════════════════════════════════════
\section{Connessioni persistenti e pipelining}
% ══════════════════════════════════════════════════════════════

Nelle connessioni \textbf{persistenti} il server \textbf{lascia aperta} la connessione TCP dopo aver
inviato una risposta.

\begin{itemize}
    \item Richieste e risposte successive tra lo stesso client e lo stesso server viaggiano sulla stessa
          connessione: l'intera pagina dell'esempio (HTML + 10 immagini) passa su \textbf{una sola}
          connessione TCP.
    \item Anche più pagine dello stesso server possono essere inviate sulla stessa connessione.
    \item L'handshake si paga \textbf{una volta sola}: ogni oggetto dopo il primo costa \textbf{1 RTT} invece di 2.
    \item Il server chiude la connessione quando resta inutilizzata per un tempo configurabile (\emph{timeout}).
\end{itemize}

\dfn{Pipelining}{
    Con il \textbf{pipelining} le richieste di più oggetti vengono inviate \textbf{una dietro l'altra},
    senza aspettare le risposte alle richieste precedenti. Da HTTP/2 richieste e risposte possono anche
    essere \textbf{intercalate} sulla stessa connessione, con un meccanismo di \textbf{priorità}. Le
    connessioni persistenti con pipelining sono oggi la modalità di default.
}

\begin{figure}[H]
\centering
\begin{tikzpicture}[yscale=0.52, every node/.style={font=\scriptsize\sffamily}]
  % non persistente
  \begin{scope}
    \node[font=\footnotesize\bfseries\sffamily] at (1.5,1) {Non persistente};
    \draw[thick] (0,0) -- (0,-15); \draw[thick] (3,0) -- (3,-15);
    \foreach \k/\t in {0/file 1, 5/file 2, 10/file 3} {
      \draw[msg, netblue] (0,-\k-0.2) -- (3,-\k-0.9); \draw[msg, netblue] (3,-\k-1.0) -- (0,-\k-1.7); \draw[msg, netblue] (0,-\k-1.8) -- (3,-\k-2.5);
      \draw[msg, netgreen, very thick] (3,-\k-2.6) -- node[above, sloped] {\t} (0,-\k-3.8);
      \node[text=netred] at (1.5,-\k-4.4) {chiusura};
    }
  \end{scope}
  % persistente
  \begin{scope}[xshift=5cm]
    \node[font=\footnotesize\bfseries\sffamily] at (1.5,1) {Persistente};
    \draw[thick] (0,0) -- (0,-15); \draw[thick] (3,0) -- (3,-15);
    \draw[msg, netblue] (0,-0.2) -- (3,-0.9); \draw[msg, netblue] (3,-1.0) -- (0,-1.7); \draw[msg, netblue] (0,-1.8) -- (3,-2.5);
    \draw[msg, netgreen, very thick] (3,-2.6) -- node[above, sloped] {file 1} (0,-3.8);
    \draw[msg, netblue] (0,-4.0) -- node[above, sloped] {richiesta 2} (3,-4.8);
    \draw[msg, netgreen, very thick] (3,-4.9) -- node[above, sloped] {file 2} (0,-6.1);
    \draw[msg, netblue] (0,-6.3) -- node[above, sloped] {richiesta 3} (3,-7.1);
    \draw[msg, netgreen, very thick] (3,-7.2) -- node[above, sloped] {file 3} (0,-8.4);
  \end{scope}
  % pipelining
  \begin{scope}[xshift=10cm]
    \node[font=\footnotesize\bfseries\sffamily] at (1.5,1) {Pipelining};
    \draw[thick] (0,0) -- (0,-15); \draw[thick] (3,0) -- (3,-15);
    \draw[msg, netblue] (0,-0.2) -- (3,-0.9); \draw[msg, netblue] (3,-1.0) -- (0,-1.7); \draw[msg, netblue] (0,-1.8) -- (3,-2.5);
    \draw[msg, netgreen, very thick] (3,-2.6) -- node[above, sloped] {file 1} (0,-3.8);
    \draw[msg, netred] (0,-4.0) -- (3,-4.8); \draw[msg, netred] (0,-4.3) -- (3,-5.1);
    \node[text=netred, rotate=-15] at (1.4,-3.7) {richieste 2 e 3};
    \draw[msg, netgreen, very thick] (3,-5.2) -- (0,-6.4); \draw[msg, netgreen, very thick] (3,-5.5) -- (0,-6.7);
    \node[text=netgreen!60!black] at (1.5,-7.3) {file 2 e 3};
  \end{scope}
  \draw[-{Stealth}, thick, black!50] (-0.7,0) -- node[left, rotate=90, anchor=south] {tempo} (-0.7,-15);
\end{tikzpicture}
\caption{Confronto indicativo degli RTT necessari per richiedere 3 file nelle tre modalità.}
\end{figure}

\subsection{Confronto}

\begin{table}[H]
\centering
\begin{tabular}{p{7cm}p{7cm}}
\toprule
\textbf{Persistente} & \textbf{Non persistente} \\
\midrule
\itick\ più veloce, soprattutto con il pipelining & \itick\ semplice da implementare, in perfetta sintonia con la natura stateless di HTTP \\
\itick\ richiede meno risorse (soprattutto memoria) & \icross\ richiede più risorse: per ogni connessione vanno allocati buffer e variabili TCP su client e server \\
\icross\ più complessa da implementare, specie con il pipelining & \icross\ più lenta: 1--2 RTT in più per ogni richiesta \\
\icross\ la connessione può restare aperta anche se inutilizzata: il server la chiude dopo un \emph{timeout} & \itick\ le connessioni non restano mai aperte inutilmente \\
\bottomrule
\end{tabular}
\caption{Connessioni persistenti e non persistenti a confronto.}
\end{table}

Lo scambio è quello di sempre: la connessione già aperta è veloce, ma è una risorsa tenuta occupata per
qualcuno che potrebbe non tornare più.

\subsection{La connessione che si chiude mentre la usi}

Che cosa succede se il server chiude per timeout una connessione persistente \textbf{proprio mentre} il
client la sta riusando?

\begin{itemize}
    \item Il server chiude da solo una connessione inattiva, senza chiedere nulla al client.
    \item Una richiesta scritta in quello stesso istante si perde con la connessione: nessuna risposta, e
          nessun errore da parte del protocollo.
    \item Un client che riusa una connessione inattiva deve quindi essere pronto ad \textbf{aprirne una nuova}
          e a \textbf{reinviare} la richiesta.
\end{itemize}

\warning{Ripetere deve essere innocuo}{
    Reinviare una richiesta è sicuro solo se ripeterla non fa danni. È uno dei motivi per cui ci si aspetta
    che una \texttt{GET} \textbf{non modifichi} nulla sul server: chiedere due volte la stessa pagina va bene,
    pagare due volte lo stesso ordine no.
}

% ══════════════════════════════════════════════════════════════
\section{HTTP/2 e HTTP/3}
% ══════════════════════════════════════════════════════════════

\subsection{Il blocco in testa alla fila}

Perché un oggetto grande può rallentare dieci oggetti piccoli? Prendiamo una pagina con un video pesante
in alto e molti piccoli oggetti sotto, una sola connessione persistente e un collegamento lento lungo il
percorso.

\dfn{Head-of-line blocking}{
    Il video occupa la connessione, e i piccoli oggetti, già pronti, aspettano dietro di lui: è il
    \textbf{blocco in testa alla fila} (\emph{head-of-line blocking}, HOL).
}

I browser aggirano il problema aprendo \textbf{fino a sei connessioni parallele}. Così, però, si prendono
anche una fetta più grande del collegamento più lento, perché TCP divide la banda in parti uguali fra le
\textbf{connessioni}, non fra le pagine. Eliminare quelle connessioni parallele senza pagare il blocco è lo
scopo di HTTP/2.

\subsection{HTTP/2: frame intercalati su una sola connessione}

\begin{itemize}
    \item Un sotto-livello di \textbf{framing} spezza ogni messaggio in piccoli \textbf{frame}: l'intestazione
          in uno, il corpo in altri, tutti in codifica \textbf{binaria}.
    \item I frame di risposte diverse vengono \textbf{intercalati} sulla stessa connessione e riassemblati
          dall'altra parte.
    \item Metodi, codici di stato, URL e campi di intestazione \textbf{non cambiano}: cambia solo il modo in
          cui i dati sono formattati e trasportati.
\end{itemize}

\begin{figure}[H]
\centering
\begin{tikzpicture}[every node/.style={font=\scriptsize\sffamily}, fr/.style={draw=black!50, minimum width=0.42cm, minimum height=0.42cm, inner sep=0pt}]
  \node[anchor=east, font=\footnotesize\bfseries\sffamily] at (-0.3,0.9) {HTTP/1.1};
  \foreach \i in {0,...,13} { \node[fr, fill=netred!35] at (\i*0.45,0.9) {}; }
  \node at (6.6,0.9) {\dots\ 1000 frame del video};
  \foreach \i in {0,...,3} { \node[fr, fill=lnet] at (8.0+\i*0.45,0.9) {}; }
  \node at (10.2,0.9) {\dots};
  \node[anchor=east, font=\footnotesize\bfseries\sffamily] at (-0.3,0) {HTTP/2};
  \foreach \i/\c in {0/netred!35,1/lnet,2/lnet,3/netred!35,4/lnet,5/lnet,6/netred!35,7/lnet,8/lnet,9/netred!35,10/lnet,11/lnet,12/netred!35,13/lnet} {
     \node[fr, fill=\c] at (\i*0.45,0) {}; }
  \node at (7.2,0) {\dots\ gli oggetti piccoli finiscono presto};
  \node[fr, fill=netred!35, label=right:{frame del video}] at (0.2,-0.8) {};
  \node[fr, fill=lnet, label=right:{frame degli oggetti piccoli}] at (3.6,-0.8) {};
\end{tikzpicture}
\caption{Con HTTP/1.1 gli oggetti piccoli aspettano la fine del video; con HTTP/2 i frame si alternano.}
\end{figure}

\begin{esempio}{18 frame invece di 1016}
Un video di 1000 frame e otto oggetti piccoli di 2 frame ciascuno. Con HTTP/1.1 gli oggetti piccoli arrivano
dopo il video: l'ultimo è completo dopo $1000 + 16 = 1016$ frame. Con HTTP/2, alternando un frame del video e
quelli degli oggetti piccoli, tutti gli otto oggetti sono consegnati dopo appena \textbf{18} frame
(i loro 16, più i 2 del video intercalati).
\end{esempio}

\subsection{Priorità e server push}

\begin{itemize}
    \item Il client può dare a ogni richiesta un \textbf{peso} fra 1 e 256 e dichiarare da quale altro
          messaggio dipende.
    \item Il server invia per primi i frame delle risposte a priorità più alta: il foglio di stile prima
          dell'immagine in fondo alla pagina.
    \item Con il \textbf{server push} il server legge l'HTML di base, capisce quali oggetti serviranno e li
          invia \textbf{prima che vengano richiesti}: si risparmia un RTT per ogni oggetto ``spinto'', cioè la
          richiesta che il client avrebbe comunque mandato.
\end{itemize}

\subsection{HTTP/3 e QUIC}

\dfn{QUIC}{
    \textbf{QUIC} è un protocollo di trasporto realizzato \textbf{nel livello applicativo}, sopra il
    semplicissimo UDP. Fornisce ciò che serve a HTTP e che UDP non dà: affidabilità, multiplazione di più
    flussi, controllo di flusso per ciascun flusso. \textbf{HTTP/3} è HTTP sopra QUIC.
}

\begin{itemize}
    \item Aprire una connessione richiede \textbf{meno RTT}, perché connessione e cifratura vengono negoziate
          \textbf{insieme} invece che una dopo l'altra.
    \item Un pacchetto perso blocca solo il \textbf{proprio} flusso: sparisce il blocco in testa alla fila
          che sopravviveva dentro TCP (in HTTP/2, una perdita ferma tutti i frame dietro di sé).
    \item Diversi meccanismi di HTTP/2 sono già dentro QUIC, quindi il protocollo sopra diventa più semplice.
\end{itemize}

\begin{table}[H]
\centering
\renewcommand{\arraystretch}{1.3}
\begin{tabular}{lll}
\toprule
\textbf{Versione} & \textbf{Connessioni per pagina} & \textbf{Che cosa ha introdotto} \\
\midrule
HTTP/1.0 (1996) & una per oggetto & connessioni non persistenti \\
HTTP/1.1 (1997) & una, riusata & connessioni persistenti, pipelining \\
HTTP/2 (2015) & una, con frame intercalati & framing binario, priorità, server push \\
HTTP/3 (2022) & una, su QUIC sopra UDP & apertura in meno RTT, niente HOL fra flussi \\
\bottomrule
\end{tabular}
\caption{Trent'anni di lavoro su un'unica domanda: quante connessioni, e per quanto tempo.}
\end{table}

% ══════════════════════════════════════════════════════════════
\section{Laboratorio: il primo GET con Wireshark}
% ══════════════════════════════════════════════════════════════

\subsection{Installare Wireshark}

Su Debian, Ubuntu e derivate (rispondere \emph{sì} quando si chiede se gli utenti non root possono catturare
pacchetti):

\begin{lstlisting}[language=bash]
$ sudo apt update && sudo apt install wireshark
$ sudo dpkg-reconfigure wireshark-common   # se si e' risposto no
$ sudo usermod -aG wireshark $USER         # poi uscire e rientrare
$ groups                                   # deve comparire wireshark
$ getcap /usr/bin/dumpcap
cap_net_raw,cap_net_admin=eip              # cio' che deve stampare
\end{lstlisting}

Se \texttt{getcap} non stampa nulla, bisogna assegnare a mano quelle due \emph{capability} al programma, con
\texttt{setcap}.

\subsection{Catturare un GET}

\begin{enumerate}
    \item Avviare Wireshark, senza far partire la cattura.
    \item Scrivere \texttt{http} nel filtro di visualizzazione, così da vedere solo i pacchetti HTTP.
    \item Aspettare poco più di un minuto (vedremo perché), poi avviare la cattura.
    \item Scaricare una pagina molto semplice: \texttt{\$ lynx http://www.columbia.edu/\textasciitilde fdc/sample.html}
    \item Fermare la cattura. Chiudere i dettagli Frame, Ethernet, IP e TCP, e aprire quello HTTP.
\end{enumerate}

Si vedono \textbf{due} pacchetti HTTP: la \texttt{GET} inviata dal client e il \texttt{200 OK} di risposta,
ciascuno dentro un segmento TCP, dentro un datagramma IP, dentro un frame Ethernet.

\begin{esercizio}{Domande sulla cattura}
Sulla \textbf{richiesta}: (a) quale versione di HTTP annuncia il client, e con quale risponde il server?
(b) Quali lingue il client dichiara di accettare? (c) Qual è l'indirizzo IP della propria macchina, e quello del
server?

Sulla \textbf{risposta}: (d) quale codice di stato ha restituito il server? (e) Quando è stata modificata
l'ultima volta la pagina? (f) Quanti byte di contenuto sono stati inviati? (g) Nei byte grezzi del pacchetto
ci sono righe di intestazione che l'elenco dei pacchetti non mostra?

\svolgimento
Le risposte vanno lette dalla \textbf{propria} cattura, non da qui. Dove guardare: la versione è nella prima
riga di ciascun messaggio (riga di richiesta e riga di stato); le lingue nella riga \texttt{Accept-Language};
gli indirizzi nel dettaglio IP (\emph{Source} e \emph{Destination}); il codice nella riga di stato; la data di
modifica in \texttt{Last-Modified}; i byte in \texttt{Content-Length}.

Per (e): la pagina può sembrare modificata meno di un minuto prima. Quel server aggiorna la data ogni minuto,
così che tutti scarichino una copia nuova: è il motivo dell'attesa al passo 3. Si ripeta poi con
\texttt{http://scottaaronson.com} e con un browser moderno al posto di \texttt{lynx}: che cosa cambia, e perché?
(Un browser moderno usa quasi sempre HTTPS: si veda il laboratorio del capitolo successivo.)
\end{esercizio}

% ══════════════════════════════════════════════════════════════
\section{Riepilogo}
% ══════════════════════════════════════════════════════════════

\riepilogo{
\begin{itemize}
    \item Il Web è una collezione di risorse identificate da \textbf{URL} e accessibili con \textbf{HTTP}, un
          protocollo client-server.
    \item HTTP usa \textbf{TCP} (porta 80) e, con HTTP/3, anche \textbf{UDP} con QUIC.
    \item HTTP è \textbf{stateless}: il server non ricorda le richieste precedenti.
    \item Connessioni \textbf{non persistenti}: una connessione per oggetto, circa 2 RTT + trasmissione per
          ciascuno. Connessioni \textbf{persistenti}: l'handshake si paga una volta, poi 1 RTT per oggetto; con il
          \textbf{pipelining} le richieste partono senza aspettare le risposte.
    \item \textbf{HTTP/2} intercala i frame di più risposte su una connessione (niente blocco in testa alla
          fila fra oggetti), con priorità e server push; \textbf{HTTP/3} porta tutto su \textbf{QUIC}, sopra UDP.
\end{itemize}
}
